Read in Gauge Theory Gravity: the metric $g$ below is the derived metric of the gravitational field $(\hb, \Omega)$ on flat spacetime, the manifold is the invertible region, and the tractor connection is the flat-space $\Cl(2,4)$ gauge field of Theorem \claimref{synthesis.gtg-tractor-gauge-field-reproduces-tractor-connection}; the index computation below is its reading in the splitting of $g$ (Fact \claimref{synthesis.imported-gtg-dictionary}).

\emph{Setting.} By Fact \claimref{interior.imported-tractor-connection}, on flat spacetime the gauge field $\Omega_{T}$ of Theorem \claimref{synthesis.gtg-tractor-gauge-field-reproduces-tractor-connection}: for a tractor $V = \sigma\nInf + \mu + \rho\nO$ and $A = \hb^{-1}(a)$, with $(u\wedge w)\cdot z = u\,(w\cdot z) - w\,(u\cdot z)$, $\ip{\nO}{\nInf} = 1$ and $\Omega(a)$ commuting with $\nO$, $\nInf$,
\[
\nabla_{A}V = (a\cdot\nabla\sigma - A\cdot\mu)\,\nInf + \bigl(D_{a}\mu + \rho A - \sigma\Sch(A)\bigr) + \bigl(a\cdot\nabla\rho + \Sch(A)\cdot\mu\bigr)\,\nO ,
\]
in components the Fact's $\nabla_{v}(\sigma,\mu,\rho) = (\nabla_{v}\sigma - \mu(v),\ \nabla_{v}\mu + v^{\flat}\rho - \Sch(v,\cdot)\sigma,\ \nabla_{v}\rho + \Sch(v,\mu^{\sharp}))$, and the connection preserves the $\Cl(2,4)$ product, the tractor form $\sigma\rho' + \rho\sigma' + \mu\cdot\mu'$, being a bivector action. Hence $\nabla_{v}X = v$ and $\nabla_{v}I_{g} = -\Sch(v)$. The classes and the Fock criterion are those of Theorem \claimref{synthesis.einstein-condition-is-infinity-vacuum-incidence}(2): a Fock vacuum $u_{W}$ is annihilated by $c(w)$ exactly for $w \in W$ (Fact \claimref{one-relation.imported-pure-spinors-incidence-chart}).

\emph{(1)} For $L = \ell$ along $\gamma$: $\nabla_{\ell}L = -Q(\ell)\,\nInf + (\ell\cdot D)\ell + (\ell\cdot\Sch(\ell))\,\nO = (\ell\cdot\Sch(\ell))\,X$, since $\ell$ is null and its physical velocity is $D$-parallel along the affinely parametrised free-fall worldline $\gamma$, and $\nabla_{\ell}X = L$. Let $u(t)$ be the transport of $u$ along $\gamma$ in the Fock module, the spin lift of the tractor transport, so that $\frac{d}{dt}\big(c(Y)u\big) = c(\nabla_{\ell}Y)u$ for every tractor field $Y$. Then $f_{1} = c(X)u$ and $f_{2} = c(L)u$ satisfy $\dot f_{1} = f_{2}$, $\dot f_{2} = (\ell\cdot\Sch(\ell))f_{1}$, a linear system with zero initial data, so $f_{1} = f_{2} = 0$ for all $t$. A conformal spinor killed by $c(X)$ and $c(L)$ has $X$ and $L$ in its three-dimensional annihilator $W_{u(t)}$, so it is a point vacuum of $\gamma(t)$, and the real null direction of $W_{u(t)}$, unique by Theorem \claimref{synthesis.einstein-condition-is-infinity-vacuum-incidence}(1), is that of $L$, namely $\gamma'(t)$. If $g$ is Einstein, $\Sch(a) = \tfrac{\Lambda}{6}a$, so $\nabla_{\ell}I_{g} = -\tfrac{\Lambda}{6}L$ and $\nabla_{\ell}L = \tfrac{\Lambda}{6}Q(\ell)X = 0$; the same argument with $f_{1} = c(I_{g})u$, $\dot f_{1} = -\tfrac{\Lambda}{6}f_{2}$, $\dot f_{2} = 0$ shows that an infinity vacuum stays one. Conversely, an infinity vacuum stays in $\ker c(I_{g})$ to first order exactly when $c(\nabla_{\ell}I_{g})u = 0$, which by Theorem \claimref{synthesis.einstein-condition-is-infinity-vacuum-incidence}(2) is the vanishing of its classes along $\ell$; for all infinity vacua at $x$ this is clause (4)(b) there, so $g$ is Einstein at $x$.

\emph{(2)} The tractor bundle, its form and its connection depend only on $[g]$ (Fact \claimref{interior.imported-tractor-connection}), and so do the Clifford and Fock modules of its fibres and their CAR, which are built from the form alone; $X$ is the canonical tractor and the point vacua are defined from $X$. The pure-scale tractor is defined by the scale, $I_{g} = \nInf$ in the splitting of $g$. For the dilation field $\hb^{-1}(b) = f^{1/2}b$ with $f = 1 + \tfrac12 x^{T}Hx$, $\hb = 1$ and $\nabla\hb = 0$ at the origin, so there $\Omega = 0$ and the field strength, hence the Schouten map, is linear in $H$. In components, a labelled reading for $\hat g = f\eta$: the Christoffel symbols vanish at the origin, $\partial_{d}\Gamma^{a}{}_{bc} = \tfrac12\eta^{aa}(H_{db}\eta_{ac} + H_{dc}\eta_{ab} - H_{da}\eta_{bc})$ and $\Ric^{\circ}_{bd} = \partial_{a}\Gamma^{a}{}_{db} - \partial_{d}\Gamma^{a}{}_{ab}$ (the convention with positive Ricci curvature for the round sphere), the store's $\Ric = -\Ric^{\circ}$ and $\Sch = \tfrac12(\Ric - \tfrac{R}{6}g)$. For $H = 2a\,e_{1}\otimes e_{1}$ (that is $f = 1 + a x_{1}^{2}$) this gives $\Sch(v) = a\,(v\cdot e^{1})\,e^{1}$ (in components $a\,dx^{1}\otimes dx^{1}$), $J = -a$, trace-free part nonzero, and $\ell\cdot\Sch(\ell) = a$ for $\ell = e_{0} + e_{1}$. For $H = -k\eta$ (the second-order germ of the constant-curvature factor $(1 + \tfrac{k}{4}\eta(x,x))^{-2}$) it gives $\Sch(v) = -\tfrac{k}{2}v$, Einstein at the origin. Since $\nabla f(0) = 0$, the change of splitting from $\eta$ to $\hat g$ is the identity at the origin, so $X$, $I_{g}$ and every vacuum there are the same for both scales; the classes differ because $\nabla_{v}I_{\hat g}(0) = -\Sch_{\hat g}(v) \ne 0 = \nabla_{v}I_{\eta}(0)$. The vacuum $u$ of $\mathrm{span}(I_{g}, L, m)$ has $c(\nabla_{\ell}I_{\eta})u = 0$ and $c(\nabla_{\ell}I_{\hat g})u \neq 0$. Any condition formulated only in terms of the conformally invariant data takes the same value for $\eta$ and $\hat g$, so it cannot imply the incidence, which holds for $\eta$ and fails for $\hat g$ at the origin.

\emph{(3)} This is the almost-Einstein characterisation (Bailey, Eastwood and Gover 1994; Gover; cited from memory), in the form of Fact \claimref{correspondence.imported-vacuum-equations-scale-tractor}, plus an algebraic identity. $\ip{I'}{X} = 1$ and $Q(I') = 2\cdot 1\cdot\tfrac{J}{4} = \tfrac{J}{2}$, so $I' - \tfrac12 Q(I')X = I_{g}$; for any $v$ with $\ip{v}{X} = 1$, $Q(v - \tfrac12 Q(v)X) = Q(v) - Q(v) = 0$, the identity behind the point embedding. Then $\nabla_{v}I' = \nabla_{v}I_{g} + \tfrac{J}{4}\nabla_{v}X + \tfrac14(v\cdot\hat\nabla J)X = -\Sch(v) + \tfrac{J}{4}v + \tfrac14(v\cdot\hat\nabla J)\,\nO$, whose middle slot is $-\Sch^{\circ}(v)$, and $\ip{\nabla_{\ell}I'}{L} = -\ell\cdot\Sch(\ell) + \tfrac{J}{4}Q(\ell) = \omega_{\ell}(I_{g},\ell)$ for null $\ell$. By Theorem \claimref{synthesis.einstein-condition-is-infinity-vacuum-incidence}(4) the incidence at $x$ is $\Sch^{\circ} = 0$ at $x$, i.e. the vanishing of the middle slot. On a connected region the same clause makes $J$ constant, so the last slot vanishes too and $I'$ is parallel; conversely a parallel section of the form $\nInf + \rho\nO$ has $\rho$ constant and $g$ Einstein with $\Lambda = 6\rho$ (Fact \claimref{correspondence.imported-vacuum-equations-scale-tractor}), which is the frozen point of Fact \claimref{mass.imported-frozen-point-tractor-connection} with $\rho_{I} = \Lambda/6$.

\emph{(4)} The formula for $\nabla_{v}I'$ is (3); it is the standard formula for the scale tractor of the unit scale (in components $\nabla_{a}I = (0, -\Sch^{\circ}_{ab}, \tfrac{1}{n}\nabla_{a}J)$, in the store's sign of the Schouten tensor), and what follows substitutes Einstein's equations into it. With $\mathcal{R}(a) = (\Lambda + \kappa\tau)a - 2\kappa T(a)$, $\kappa = 4\pi G$: the scalar curvature is $4\Lambda + 2\kappa\tau$, $\Sch(a) = \tfrac12(\mathcal{R}(a) - \tfrac16(4\Lambda + 2\kappa\tau)a)$, so $\Sch^{\circ}$, half the trace-free part of the Ricci map, is $-\kappa T^{\circ}$, and $J = \tfrac16(4\Lambda + 2\kappa\tau)$, $\tfrac{J}{4} = \tfrac{\Lambda}{6} + \tfrac{\kappa\tau}{12}$. Hence the middle slot of $\nabla_{v}I'$ is $\kappa T^{\circ}(v)$, the last slot is $\tfrac{\kappa}{12}\,v\cdot\hat\nabla\tau$ (as $\Lambda$ is constant by clause (7) there), and $6\rho_{I'} - \tfrac{\kappa}{2}\tau = \Lambda$. Conversely $-\Sch^{\circ} = \kappa T^{\circ}$ says $\mathcal{R}(a) + 2\kappa T(a)$ is a multiple of $a$, the pointwise condition of clause (7). For null $\ell$, $\ell\cdot T^{\circ}(\ell) = \ell\cdot T(\ell)$, so $\ip{\nabla_{\ell}I'}{L} = \kappa\,\ell\cdot T(\ell)$.

\emph{The check.} Exact Gaussian-rational arithmetic and GA exterior-model Fock states. (1) For eleven rational null $\ell$ and a generic and an Einstein Schouten tensor: $\omega(\ell)X = L$, $\omega(\ell)L = \Sch(\ell,\ell)X$ and closure of $\mathrm{span}(X, L)$; closure of $\mathrm{span}(I_{g}, L)$ for the Einstein tensor and its failure for the generic one; and in the Fock module, the point vacuum of $\mathrm{span}(X, L, m)$ killed by both derivatives for both tensors, the infinity vacuum of $\mathrm{span}(I_{g}, L, m)$ killed only for the Einstein one. (2) The Ricci tensor at the origin of $f\eta$ from the Hessian; trace-free Schouten zero for $H = -k\eta$ at two values of $k$ and nonzero for $f = 1 + 3x_{1}^{2}$, with own-ray class nonzero there and zero for $\eta$, and the same Fock vacuum killed by $c(\nabla_{\ell}I_{\eta})$ but not by $c(\nabla_{\ell}I_{\hat g})$. (3) $\ip{I'}{X} = 1$, $Q(I') = J/2$, $I_{g} = I' - \tfrac12 Q(I')X$, the point-embedding identity on a sample tractor, equality of own-ray classes of $I'$ and $I_{g}$ on eleven null rays, the middle slot $-\Sch^{\circ}(v,\cdot)$ on three directions, and $\nabla I' = 0$ for the Einstein tensor with $\Lambda = 3$. (4) For two rational pairs $(\kappa, \Lambda)$ and a rational $T$: the middle slot $\kappa T^{\circ}$, $J/4 = \Lambda/6 + \kappa\tau/12$, $6\rho_{I'} - \kappa\tau/2 = \Lambda$, and the own-ray classes $\kappa T(\ell,\ell)$.
